\chapter{Przestrzenie topologiczne}
\label{ch:topog-przestrzenie-topologiczne}

\noindent\textbf{Cel rozdziału.}
W analizie pojęcia granicy i ciągłości poznajemy najpierw za pomocą odległości.
Zaczynamy od elementarnego wprowadzenia do przestrzeni metrycznych:
metryki, kul i sfer oraz przykładów na prostej i w przestrzeni euklidesowej.
Następnie sprawdzamy, że wiele argumentów nie wymaga liczbowej
wartości odległości: wystarcza wiedzieć, które zbiory są otwarte.
Precyzujemy aksjomaty topologii, wyprowadzamy pierwsze konsekwencje,
konstruujemy topologie z baz i podbaz oraz porównujemy różne topologie
na tym samym zbiorze. Zaawansowane własności przestrzeni metrycznych
(zbieżność, zupełność, metryzowalność) omówimy w dalszych rozdziałach.
Nie zakładamy żadnych aksjomatów oddzielania, przeliczalności
ani niepustości przestrzeni.

\section{Przestrzenie metryczne: definicje i podstawowe przykłady}
\label{sec:topog-motywacja}

Zanim zdefiniujemy topologię, wprowadzimy znany z geometrii
sposób mierzenia odległości. Posłuży nam on jako punkt wyjścia,
ale nie będziemy zakładać istnienia odległości w dowolnej
przestrzeni topologicznej.

\begin{definicja}[Metryka i przestrzeń metryczna]
\label{def:topog-metryka}
Niech $X$ będzie zbiorem. \emph{Metryką} na $X$ nazywamy
odwzorowanie $d\colon X\times X\to[0,\infty)$ takie, że
dla wszystkich $x,y,z\in X$ zachodzą:
\begin{enumerate}[label=(M\arabic*)]
\item $d(x,y)=0$ wtedy i tylko wtedy, gdy $x=y$;
\item $d(x,y)=d(y,x)$ (symetria);
\item $d(x,z)\le d(x,y)+d(y,z)$ (nierówność trójkąta).
\end{enumerate}
Parę $(X,d)$ nazywamy \emph{przestrzenią metryczną}.
Liczba $d(x,y)$ jest \emph{odległością} między $x$ a $y$.
Dopuszczamy $X=\varnothing$; jedyna funkcja
$\varnothing\times\varnothing\to[0,\infty)$ jest wtedy metryką.
\end{definicja}

\begin{przyklad}[Prosta rzeczywista i przestrzeń euklidesowa]
\label{ex:topog-metryka-euklidesowa}
Na $\R$ funkcja $d(x,y)=|x-y|$ jest metryką.
Warunek (M1) wynika z własności wartości bezwzględnej,
(M2) z równości $|x-y|=|y-x|$, a (M3) z nierówności
$|x-z|\le |x-y|+|y-z|$.

Dla $n\ge1$ na $\R^n$ kładziemy
$d_2(x,y)=\sqrt{\sum_{j=1}^{n}(x_j-y_j)^2}=\|x-y\|_2$.
Symetria i warunek (M1) są natychmiastowe.
Uzasadnijmy używaną tu nierówność Cauchy'ego--Schwarza.
Dla $v\ne0$, przyjmując
$t=\bigl(\sum_j u_jv_j\bigr)/\sum_j v_j^2$, z nierówności
$0\le\sum_j(u_j-tv_j)^2$ otrzymujemy
$(\sum_j u_jv_j)^2\le(\sum_j u_j^2)(\sum_j v_j^2)$;
dla $v=0$ nierówność jest oczywista.
Zatem $\|u+v\|_2^2=\|u\|_2^2+2\sum_j u_jv_j+\|v\|_2^2
\le(\|u\|_2+\|v\|_2)^2$.
Po wzięciu pierwiastków i podstawieniu
$u=x-y$, $v=y-z$ otrzymujemy (M3).
Dla $n=1$ wzór ten sprowadza się do wartości bezwzględnej.
\end{przyklad}

\begin{przyklad}[Metryka maksimum]
\label{ex:topog-metryka-maksimum}
Na $\R^n$, $n\ge1$, określmy
$d_\infty(x,y)=\max_{1\le j\le n}|x_j-y_j|$.
Warunki (M1) i (M2) wynikają z własności wartości
bezwzględnej. Dla każdych $x,y,z\in\R^n$ i każdego $j$
mamy $|x_j-z_j|\le|x_j-y_j|+|y_j-z_j|
\le d_\infty(x,y)+d_\infty(y,z)$.
Biorąc maksimum po $j$, dostajemy (M3).
W szczególności na $\R^2$ kula tej metryki ma kształt
kwadratu o bokach równoległych do osi współrzędnych.
\end{przyklad}

\begin{przyklad}[Metryka dyskretna]
\label{ex:topog-metryka-dyskretna}
Na dowolnym zbiorze $X$ określmy
$d_{\mathrm{dys}}(x,y)=0$, gdy $x=y$, oraz
$d_{\mathrm{dys}}(x,y)=1$, gdy $x\ne y$.
Symetria i (M1) są oczywiste. Jeśli $x\ne z$, to
przynajmniej jedna z nierówności $x\ne y$ lub $y\ne z$
jest prawdziwa. Stąd
$1=d_{\mathrm{dys}}(x,z)\le
 d_{\mathrm{dys}}(x,y)+d_{\mathrm{dys}}(y,z)$.
Dla $x=z$ nierówność trójkąta także zachodzi.
Metryka dyskretna nie wymaga żadnej struktury algebraicznej
ani porządku na $X$.
\end{przyklad}

\begin{definicja}[Kule i sfera metryczna]
\label{def:topog-kule-sfera}
Niech $(X,d)$ będzie przestrzenią metryczną,
$x\in X$ i $r>0$. Definiujemy:
\begin{align*}
B_d(x,r)&=\{y\in X:d(x,y)<r\},
&&\text{\emph{kula otwarta}},\\
\overline B_d(x,r)&=\{y\in X:d(x,y)\le r\},
&&\text{\emph{kula domknięta}},\\
S_d(x,r)&=\{y\in X:d(x,y)=r\},
&&\text{\emph{sfera metryczna}}.
\end{align*}
Punkt $x$ nazywamy \emph{środkiem}, a $r$ \emph{promieniem}.
W każdym z tych zbiorów bierzemy pod uwagę wyłącznie
punkty należące do $X$.
\end{definicja}

\begin{przyklad}[Kule na prostej i płaszczyźnie]
\label{ex:topog-kule-euklidesowe}
Dla $X=\R$ i $d(x,y)=|x-y|$ mamy
$B_d(x,r)=(x-r,x+r)$,
$\overline B_d(x,r)=[x-r,x+r]$ oraz
$S_d(x,r)=\{x-r,x+r\}$.
Na płaszczyźnie euklidesowej $B_{d_2}(x,r)$ jest
wnętrzem dysku o środku $x$ i promieniu $r$,
$\overline B_{d_2}(x,r)$ jest całym dyskiem,
a $S_{d_2}(x,r)$ jest okręgiem.
W $\R^3$ sfera $S_{d_2}(x,r)$ jest powierzchnią kuli,
nie jej wnętrzem. Dla metryki maksimum na $\R^2$
kula $B_{d_\infty}(x,r)$ jest otwartym kwadratem,
a $S_{d_\infty}(x,r)$ jego brzegiem geometrycznym.
\end{przyklad}

\begin{figure}[ht]
\centering
\begin{tikzpicture}[>=Latex, font=\small]
  \begin{scope}[shift={(0,0)}]
    \fill[blue!10] (0,0) circle (0.83);
    \draw[blue!65!black,densely dashed,thick] (0,0) circle (0.83);
    \fill (0,0) circle (1.2pt);
    \draw[->,thin] (0,0)--(.58,.58) node[pos=.61,above left] {$r$};
    \node at (0,-1.20) {$B_d(x,r)$};
    \node[font=\scriptsize,align=center] at (0,-1.60) {bez okręgu};
  \end{scope}
  \begin{scope}[shift={(3.4,0)}]
    \fill[blue!10] (0,0) circle (0.83);
    \draw[blue!65!black,thick] (0,0) circle (0.83);
    \fill (0,0) circle (1.2pt);
    \draw[->,thin] (0,0)--(.58,.58);
    \node at (0,-1.20) {$\overline B_d(x,r)$};
    \node[font=\scriptsize] at (0,-1.60) {z okręgiem};
  \end{scope}
  \begin{scope}[shift={(6.8,0)}]
    \draw[blue!65!black,thick] (0,0) circle (0.83);
    % Środek służy tylko jako odniesienie: nie jest punktem sfery.
    \draw[gray!75,fill=white] (0,0) circle (1.2pt);
    \draw[->,thin] (0,0)--(.58,.58);
    \node at (0,-1.20) {$S_d(x,r)$};
    \node[font=\scriptsize] at (0,-1.60) {tylko okrąg};
  \end{scope}
\end{tikzpicture}
\caption{Kula otwarta, kula domknięta i sfera
w płaszczyźnie z metryką euklidesową. Na rysunku
linia przerywana nie należy do kuli otwartej, a zaznaczony
środek nie należy do sfery. Dla ogólnej metryki kształt tych
zbiorów może być inny.}
\label{fig:topog-kule-sfera}
\end{figure}

\begin{definicja}[Zbiór otwarty w przestrzeni metrycznej]
\label{def:topog-otwarty-metrycznie}
Niech $(X,d)$ będzie przestrzenią metryczną.
Podzbiór $U\subseteq X$ nazywamy \emph{otwartym
w sensie metrycznym}, jeśli dla każdego $x\in U$
istnieje $r>0$ takie, że $B_d(x,r)\subseteq U$.
W szczególności $\varnothing$ jest otwarty:
nie ma w nim punktów wymagających sprawdzenia.
Podzbiór $F\subseteq X$ nazywamy \emph{domkniętym
w sensie metrycznym}, jeśli $X\setminus F$ jest otwarty.
\end{definicja}

\begin{stwierdzenie}[Otwartość i domkniętość kul oraz sfer]
\label{prop:topog-kule-otwarte-sfery-domkniete}
Niech $(X,d)$ będzie przestrzenią metryczną,
$x\in X$ i $r>0$. Wówczas kula $B_d(x,r)$ jest otwarta
w sensie metrycznym, a zbiory $\overline B_d(x,r)$
i $S_d(x,r)$ są domknięte w sensie metrycznym.
\end{stwierdzenie}
\begin{proof}
Z nierówności trójkąta i symetrii otrzymujemy dla
wszystkich $y,z\in X$
$|d(x,y)-d(x,z)|\le d(y,z)$.
Jeżeli $y\in B_d(x,r)$, to $\varepsilon=r-d(x,y)>0$.
Dla każdego $z\in B_d(y,\varepsilon)$ zachodzi
$d(x,z)\le d(x,y)+d(y,z)<r$, więc
$B_d(y,\varepsilon)\subseteq B_d(x,r)$.

Jeżeli $y\notin\overline B_d(x,r)$, to
$\varepsilon=d(x,y)-r>0$. Dla każdego
$z\in B_d(y,\varepsilon)$ mamy
$d(x,z)\ge d(x,y)-d(y,z)>r$.
Dopełnienie kuli domkniętej jest zatem otwarte.

Wreszcie, jeśli $y\notin S_d(x,r)$, to
$\varepsilon=|d(x,y)-r|>0$. Każdy
$z\in B_d(y,\varepsilon)$ spełnia
$|d(x,z)-d(x,y)|<\varepsilon$,
więc $d(x,z)$ leży po tej samej stronie $r$ co $d(x,y)$.
W szczególności $z\notin S_d(x,r)$.
Zatem dopełnienie sfery jest otwarte.
\end{proof}

\begin{uwaga}[Kula domknięta nie musi być domknięciem kuli otwartej]
\label{rem:topog-kula-domknieta-a-domkniecie}
Nazwy „kula domknięta” i „sfera” odnoszą się do
warunków określających odległość od środka. W ogólnej przestrzeni
metrycznej $\overline B_d(x,r)$ nie musi być domknięciem
$B_d(x,r)$, a sfera nie musi być brzegiem kuli otwartej.
Przykładowo, dla metryki dyskretnej na zbiorze $X$
mającym co najmniej dwa punkty i promienia $r=1$ mamy
$B_d(x,1)=\{x\}$, $\overline B_d(x,1)=X$ i
$S_d(x,1)=X\setminus\{x\}$.
Singleton $\{x\}$ jest otwarty, a także domknięty,
bo $X\setminus\{x\}$ jest sumą otwartych kul o promieniu
$1/2$ ze środkami w pozostałych punktach. Wobec tego
$\{x\}$ jest już zbiorem domkniętym, podczas gdy
kula domknięta $\overline B_d(x,1)$ jest całym $X$.
To rozróżnienie oraz pojęcie brzegu dokładniej omówimy
w rozdziale 02.
\end{uwaga}

\section{Od metryki do aksjomatów topologii}
\label{sec:topog-odleglosc-otwarte}

W przestrzeni metrycznej zadana metryka pozwala określić
zbiory otwarte. Okazuje się, że rodzina tych zbiorów
spełnia proste aksjomaty, które można przyjąć bez
odwoływania się do odległości.

\begin{stwierdzenie}[Zbiory otwarte w przestrzeni metrycznej]
\label{prop:topog-topologia-metryczna}
Niech $(X,d)$ będzie dowolną przestrzenią metryczną, a $\tau_d$
rodziną wszystkich podzbiorów $X$ otwartych w sensie
Definicji~\ref{def:topog-otwarty-metrycznie}.
Wówczas:
\begin{enumerate}[label=(\roman*)]
\item $\varnothing,X\in\tau_d$;
\item suma dowolnej rodziny zbiorów z $\tau_d$ należy do $\tau_d$;
\item przecięcie dowolnej skończonej rodziny zbiorów z $\tau_d$
      należy do $\tau_d$.
\end{enumerate}
\end{stwierdzenie}
\begin{proof}
Zbiór pusty jest otwarty na mocy definicji. Jeżeli $x\in X$,
to $B_d(x,1)\subseteq X$, więc $X$ również jest otwarty;
dotyczy to także $X=\varnothing$.

Niech $(U_i)_{i\in I}$ będzie dowolną rodziną zbiorów otwartych.
Jeżeli $x\in\bigcup_{i\in I}U_i$, to $x\in U_j$ dla pewnego $j\in I$.
Istnieje $r>0$ takie, że $B_d(x,r)\subseteq U_j\subseteq
\bigcup_{i\in I}U_i$. A zatem suma jest otwarta.
Gdy $I=\varnothing$, suma jest zbiorem pustym.

Dla dwóch zbiorów otwartych $U,V$ i punktu $x\in U\cap V$
wybieramy $r,s>0$ takie, że $B_d(x,r)\subseteq U$ i
$B_d(x,s)\subseteq V$. Wówczas
$B_d(x,\min\{r,s\})\subseteq U\cap V$.
Indukcja daje tezę dla dowolnej niepustej skończonej rodziny;
przecięciem rodziny pustej jest $X$, które jest otwarte.
\end{proof}

\begin{uwaga}
\label{rem:topog-nie-metryczna}
Aksjomaty z powyższego stwierdzenia mają sens na dowolnym zbiorze,
nawet gdy nie jest na nim zadana metryka. To prowadzi do podstawowej
definicji topologii. W dalszych rozdziałach zobaczymy także topologie,
których nie można otrzymać z żadnej metryki.
\end{uwaga}

\section{Aksjomaty przestrzeni topologicznej}
\label{sec:topog-aksjomaty}

\begin{definicja}[Topologia i przestrzeń topologiczna]
\label{def:topog-topologia}
Niech $X$ będzie dowolnym zbiorem. \emph{Topologią na $X$}
nazywamy rodzinę $\tau\subseteq\mathcal P(X)$ spełniającą:
\begin{enumerate}[label=(T\arabic*)]
\item $\varnothing\in\tau$ i $X\in\tau$;
\item dla każdej rodziny $(U_i)_{i\in I}$ elementów $\tau$
      zbiór $\bigcup_{i\in I}U_i$ należy do $\tau$;
\item dla każdych $U,V\in\tau$ zbiór $U\cap V$ należy do $\tau$.
\end{enumerate}
Parę $(X,\tau)$ nazywamy \emph{przestrzenią topologiczną},
a elementy $\tau$ \emph{zbiorami otwartymi}. Zbiór $F\subseteq X$
nazywamy \emph{domkniętym}, jeśli $X\setminus F\in\tau$.
Jeżeli topologia nie budzi wątpliwości, piszemy krótko $X$
zamiast $(X,\tau)$.
\end{definicja}

Warunek (T3) jest równoważny zamknięciu na \emph{skończone}
przecięcia: wykonujemy kolejne przecięcia dwóch zbiorów,
a przecięcie pustej rodziny to $X$. W (T2) indeks $I$
może być dowolny, także pusty. Nie wymagamy, by przecięcie
nieskończenie wielu zbiorów otwartych było otwarte.

\begin{przyklad}[Przecięcie nieskończenie wielu zbiorów otwartych]
\label{ex:topog-nieskonczone-przeciecie}
W $\R$ z topologią pochodzącą od metryki $d(x,y)=|x-y|$
zbiory $U_n=(-1/n,1/n)$, $n\geq1$, są otwarte, lecz
$\bigcap_{n\geq1}U_n=\{0\}$ nie jest otwarty.
Istotnie, dla każdego $r>0$ punkt $r/2$ należy do
$B_d(0,r)$, ale nie do $\{0\}$.
Ponadto $x\neq0$ nie należy do $U_n$ dla dostatecznie
dużego $n$, dlatego przecięcie jest dokładnie $\{0\}$.
\end{przyklad}

\begin{stwierdzenie}[Aksjomaty zbiorów domkniętych]
\label{prop:topog-zbiory-domkniete}
Niech $(X,\tau)$ będzie przestrzenią topologiczną, a
$\mathcal F=\{F\subseteq X:X\setminus F\in\tau\}$.
Wówczas $\varnothing,X\in\mathcal F$, przecięcie dowolnej
rodziny elementów $\mathcal F$ należy do $\mathcal F$
i suma dwóch elementów $\mathcal F$ należy do $\mathcal F$.
Odwrotnie, każda rodzina $\mathcal F\subseteq\mathcal P(X)$
spełniająca te trzy warunki wyznacza dokładnie jedną topologię,
której rodziną zbiorów domkniętych jest $\mathcal F$.
\end{stwierdzenie}
\begin{proof}
Pierwsze trzy własności wynikają z praw de Morgana:
$X\setminus\bigcap_{i\in I}F_i
=\bigcup_{i\in I}(X\setminus F_i)$ i
$X\setminus(F\cup G)=(X\setminus F)\cap(X\setminus G)$.
Pierwsza równość jest prawdziwa również dla pustej rodziny.

Odwrotnie, dla zadanej $\mathcal F$ definiujemy
$\tau=\{X\setminus F:F\in\mathcal F\}$. Z warunku
$\varnothing,X\in\mathcal F$ otrzymujemy (T1),
z zamknięcia $\mathcal F$ na dowolne przecięcia otrzymujemy
(T2), a z jej zamknięcia na sumy dwóch zbiorów otrzymujemy
(T3). Ponieważ $\varnothing\in\mathcal F$, rodzina
$\mathcal F$ jest także zamknięta na sumy skończone.
Rodzina zbiorów domkniętych tej topologii jest
równa $\mathcal F$. Każda topologia o tej rodzinie
zbiorów domkniętych musi składać się z ich dopełnień,
więc jest jednoznacznie wyznaczona.
\end{proof}

\section{Pierwsze przykłady i kontrprzykłady}
\label{sec:topog-przyklady}

\begin{przyklad}[Topologia dyskretna i antydyskretna]
\label{ex:topog-dyskretna-antydyskretna}
Dla dowolnego zbioru $X$ rodzina $\mathcal P(X)$ jest topologią,
ponieważ wszystkie występujące w aksjomatach zbiory są
podzbiorami $X$. Jest to \emph{topologia dyskretna}.
Jest ona topologią pochodzącą od metryki dyskretnej
z Przykładu~\ref{ex:topog-metryka-dyskretna}: dla każdego $x\in X$
mamy $B_{d_{\mathrm{dys}}}(x,1/2)=\{x\}$, a więc
każdy podzbiór $U\subseteq X$ jest otwarty, także gdy $U=\varnothing$.
Rodzina $\{\varnothing,X\}$ także spełnia (T1)--(T3):
sumy i przecięcia jej elementów należą do tej rodziny.
Jest to \emph{topologia antydyskretna}.
Jeżeli $X=\varnothing$ lub $X$ jest jednoelementowy,
obie topologie są identyczne; jeżeli $X$ ma co najmniej
dwa elementy, są różne.
\end{przyklad}

\begin{przyklad}[Topologia koskończona]
\label{ex:topog-kofinita}
Niech $X$ będzie dowolnym zbiorem. Ustalmy
\[
\tau_{\mathrm{ks}}=\{\varnothing\}\cup
\{U\subseteq X:X\setminus U\text{ jest skończony}\}.
\]
Zbiory $\varnothing$ i $X$ należą do $\tau_{\mathrm{ks}}$.
Jeżeli rodzina elementów $\tau_{\mathrm{ks}}$ zawiera
niepusty składnik $U$, to dopełnienie jej sumy jest zawarte
w skończonym zbiorze $X\setminus U$, a zatem jest skończone.
Jeżeli wszystkie składniki są puste albo rodzina jest pusta,
ich suma jest pusta.
Dla niepustych $U,V\in\tau_{\mathrm{ks}}$ dopełnienie
$U\cap V$ jest sumą dwóch zbiorów skończonych.
Przecięcie z $\varnothing$ również jest otwarte.
Są to dokładnie aksjomaty topologii.

Gdy $X$ jest skończony, $\tau_{\mathrm{ks}}=\mathcal P(X)$.
Gdy $X$ jest nieskończony, dwa niepuste zbiory otwarte
nie mogą być rozłączne: ich dopełnienia są skończone,
a zbiory dopełniające nie mogłyby razem pokryć
nieskończonego $X$.
\end{przyklad}

\begin{przyklad}[Topologia koprzeliczalna]
\label{ex:topog-koprzeliczalna}
Niech $X$ będzie dowolnym zbiorem. Definiujemy
\[
\tau_{\mathrm{kp}}=\{\varnothing\}\cup
\{X\setminus C:C\subseteq X,\ C\text{ jest co najwyżej przeliczalny}\}.
\]
Argument z poprzedniego przykładu działa po zastąpieniu
słowa „skończony” słowami „co najwyżej przeliczalny”:
podzbiór zbioru co najwyżej przeliczalnego jest
co najwyżej przeliczalny, a suma dwóch zbiorów
co najwyżej przeliczalnych jest co najwyżej przeliczalna.
Zatem jest to topologia. Jeśli $X$ jest co najwyżej
przeliczalny, otrzymujemy topologię dyskretną.
Jeżeli $X$ jest nieprzeliczalny, to każde dwa niepuste
zbiory otwarte mają niepuste przecięcie, ponieważ suma
dwóch co najwyżej przeliczalnych dopełnień nie może pokryć $X$.
\end{przyklad}

\begin{przyklad}[Przestrzeń Sierpińskiego]
\label{ex:topog-sierpinski}
Na $X=\{0,1\}$ określmy
$\tau_{\mathrm{Si}}=\{\varnothing,\{1\},X\}$. Łatwo bezpośrednio
sprawdzić, że suma i przecięcie dowolnych elementów
tej rodziny należą do $\tau_{\mathrm{Si}}$: jedynymi możliwościami
są $\varnothing,\{1\},X$. Jest to topologia różna
zarówno od dyskretnej, jak i antydyskretnej.
Zbiór $\{1\}$ jest otwarty, lecz nie jest domknięty,
ponieważ jego dopełnienie $\{0\}$ nie jest otwarte.
Natomiast $\{0\}$ jest domknięty, lecz nieotwarty.
\end{przyklad}

\begin{przyklad}[Topologia wyznaczona przez podział]
\label{ex:topog-podzial}
Niech $\mathcal A$ będzie podziałem zbioru $X$:
każdy $A\in\mathcal A$ jest niepusty, różne elementy
$\mathcal A$ są rozłączne, a $\bigcup\mathcal A=X$.
Rodzina $\tau_{\mathcal A}=\{\bigcup\mathcal E:
\mathcal E\subseteq\mathcal A\}$ jest topologią.
Zbiór pusty powstaje jako suma pustej podrodziny,
a $X$ jako suma całego podziału. Suma dowolnych sum
bloków znów jest sumą bloków. Przecięcie dwóch takich
sum jest sumą bloków należących do obu wybranych
podrodzin: każdy punkt należy do dokładnie jednego bloku.
Gdy $\mathcal A=\{X\}$ i $X\neq\varnothing$, mamy
topologię antydyskretną; gdy
$\mathcal A=\{\{x\}:x\in X\}$, mamy topologię dyskretną.
\end{przyklad}

\begin{uwaga}[Przestrzeń pusta]
\label{rem:topog-pusta}
Na $X=\varnothing$ istnieje tylko jedna topologia,
$\tau=\{\varnothing\}$. W aksjomacie (T1) zbiory
$\varnothing$ i $X$ są wtedy tym samym zbiorem.
Uwzględnianie przypadku pustego zapobiega
niepotrzebnym wyjątkom w ogólnych konstrukcjach.
\end{uwaga}

\section{Otoczenia punktów}
\label{sec:topog-otoczenia}

\begin{definicja}[Otoczenie i otoczenie otwarte]
\label{def:topog-otoczenie}
Niech $(X,\tau)$ będzie przestrzenią topologiczną i $x\in X$.
Zbiór $N\subseteq X$ nazywamy \emph{otoczeniem punktu $x$},
jeśli istnieje $U\in\tau$ takie, że $x\in U\subseteq N$.
Jeżeli dodatkowo $N\in\tau$, mówimy o
\emph{otoczeniu otwartym}. Rodzinę wszystkich otoczeń
punktu $x$ oznaczamy przez $\mathcal N(x)$.
\end{definicja}

Nie każde otoczenie jest otwarte: w standardowej topologii
na $\R$ zbiór $[-1,1]$ jest otoczeniem zera, ponieważ
$0\in(-1,1)\subseteq[-1,1]$, lecz sam nie jest otwarty.

\begin{stwierdzenie}[Podstawowe własności otoczeń]
\label{prop:topog-otoczenia}
Dla dowolnej przestrzeni topologicznej $(X,\tau)$
i każdego $x\in X$ zachodzą następujące własności:
\begin{enumerate}[label=(\roman*)]
\item $X\in\mathcal N(x)$, a każde $N\in\mathcal N(x)$ zawiera $x$;
\item jeżeli $N\in\mathcal N(x)$ oraz $N\subseteq M\subseteq X$,
      to $M\in\mathcal N(x)$;
\item jeżeli $N,M\in\mathcal N(x)$, to
      $N\cap M\in\mathcal N(x)$;
\item jeżeli $N\in\mathcal N(x)$, to istnieje $M\in\mathcal N(x)$
      takie, że $N\in\mathcal N(y)$ dla każdego $y\in M$.
\end{enumerate}
Zbiór $U\subseteq X$ jest otwarty wtedy i tylko wtedy, gdy
$U\in\mathcal N(x)$ dla każdego $x\in U$.
\end{stwierdzenie}
\begin{proof}
Ponieważ $X$ jest otwarty, jest otoczeniem każdego swojego punktu.
Pozostała część (i) wynika wprost z definicji.
Jeżeli $x\in V\subseteq N\subseteq M$ i $V$ jest otwarty,
to $M$ również jest otoczeniem $x$, co daje (ii).
Dla (iii) wybieramy otwarte $V,W$ takie, że
$x\in V\subseteq N$ i $x\in W\subseteq M$.
Zbiór $V\cap W$ jest otwarty i
$x\in V\cap W\subseteq N\cap M$.
W (iv) możemy wziąć $M=V$, gdzie $V$ jest dowolnym
otwartym zbiorem spełniającym $x\in V\subseteq N$:
dla każdego $y\in V$ ten sam $V$ świadczy,
że $N$ jest otoczeniem $y$.

Jeżeli $U$ jest otwarty, jest otoczeniem każdego punktu,
który zawiera. Odwrotnie, załóżmy, że $U\in\mathcal N(x)$
dla każdego $x\in U$. Każdy taki punkt należy do pewnego
zbioru otwartego $V$ zawartego w $U$. Wobec tego
\[
U=\bigcup\{V\in\tau:V\subseteq U\},
\]
więc $U$ jest otwarty. Równość pozostaje prawdziwa także dla
$U=\varnothing$: rodzina po prawej stronie jest wtedy
jednoelementowa i składa się ze zbioru pustego.
\end{proof}

\begin{uwaga}
\label{rem:topog-otoczenia-topologia}
Własność (iv) odzwierciedla fakt, że punkt nie musi być
jedynym punktem posiadającym dane otoczenie: pewien
jego mniejszy fragment jest otoczeniem każdego swojego punktu.
Pełne odwrotne twierdzenie o opisywaniu topologii za pomocą
rodzin otoczeń jest użyteczne przy przejściu do przestrzeni
jednostajnych; na obecnym etapie wystarczy powyższa charakteryzacja.
\end{uwaga}

\section{Bazy topologii}
\label{sec:topog-bazy}

\begin{definicja}[Baza topologii]
\label{def:topog-baza}
Niech $(X,\tau)$ będzie przestrzenią topologiczną.
Rodzinę $\mathcal B\subseteq\tau$ nazywamy
\emph{bazą topologii $\tau$}, jeśli dla każdego
$U\in\tau$ oraz $x\in U$ istnieje $B\in\mathcal B$
takie, że $x\in B\subseteq U$.
\end{definicja}

Równoważnie, każdy zbiór otwarty jest sumą pewnej podrodziny
$\mathcal B$. Istotnie, wybierając wszystkie elementy bazy
zawarte w $U$, otrzymujemy sumę równą $U$;
w drugą stronę, jeśli $U$ jest taką sumą, każdy
jego punkt należy do któregoś składnika.

\begin{twierdzenie}[Kryterium bycia bazą]
\label{thm:topog-kryterium-bazy}
Niech $X$ będzie zbiorem, a $\mathcal B\subseteq\mathcal P(X)$.
Istnieje topologia na $X$, dla której $\mathcal B$ jest bazą,
wtedy i tylko wtedy, gdy:
\begin{enumerate}[label=(B\arabic*)]
\item dla każdego $x\in X$ istnieje $B\in\mathcal B$
      takie, że $x\in B$;
\item jeżeli $B_1,B_2\in\mathcal B$ oraz
      $x\in B_1\cap B_2$, to istnieje $B_3\in\mathcal B$
      takie, że $x\in B_3\subseteq B_1\cap B_2$.
\end{enumerate}
Jeżeli te warunki są spełnione, odpowiednia topologia
jest wyznaczona jednoznacznie i ma postać
\[
\tau_{\mathcal B}=
\left\{\bigcup\mathcal E:\mathcal E\subseteq\mathcal B\right\}.
\]
\end{twierdzenie}
\begin{proof}
Jeżeli $\mathcal B$ jest bazą topologii $\tau$, to dla
$U=X$ jej definicja daje (B1). Jeżeli
$x\in B_1\cap B_2$, zbiór $B_1\cap B_2$ jest otwarty,
więc z definicji bazy istnieje $B_3$ jak w (B2).

Załóżmy odwrotnie (B1)--(B2) i zdefiniujmy
$\tau_{\mathcal B}$ podanym wzorem. Suma pustej
podrodziny daje $\varnothing$; z (B1) wynika
$X=\bigcup\mathcal B\in\tau_{\mathcal B}$.
Suma dowolnej rodziny elementów $\tau_{\mathcal B}$
jest sumą wszystkich występujących w ich zapisie
elementów bazy, więc należy do $\tau_{\mathcal B}$.

Niech $U=\bigcup\mathcal E$ oraz $V=\bigcup\mathcal D$,
gdzie $\mathcal E,\mathcal D\subseteq\mathcal B$.
Dla każdego $x\in U\cap V$ istnieją
$B_1\in\mathcal E$ i $B_2\in\mathcal D$
takie, że $x\in B_1\cap B_2$. Z (B2) wynika więc
istnienie $B\in\mathcal B$ spełniającego
$x\in B\subseteq B_1\cap B_2\subseteq U\cap V$.
W konsekwencji
\[
U\cap V=\bigcup\{B\in\mathcal B:B\subseteq U\cap V\},
\]
co dowodzi, że $U\cap V\in\tau_{\mathcal B}$.
Ten zapis obejmuje również przypadek $U\cap V=\varnothing$
i nie wymaga wyboru po jednym elemencie bazy dla każdego punktu.
Otrzymujemy topologię. Wszystkie elementy
$\mathcal B$ są w niej otwarte, a każdy zbiór otwarty
jest sumą elementów $\mathcal B$, więc $\mathcal B$
jest jej bazą. Jednoznaczność wynika z równoważnej
charakteryzacji bazy podanej przed twierdzeniem.
\end{proof}

\begin{stwierdzenie}[Kryterium otwartości przez bazę]
\label{prop:topog-otwartosc-przez-baze}
Niech $\mathcal B$ będzie bazą topologii $\tau$ na $X$.
Dla podzbioru $U\subseteq X$ następujące warunki są równoważne:
\begin{enumerate}[label=(\roman*)]
\item $U\in\tau$;
\item dla każdego $x\in U$ istnieje $B\in\mathcal B$
      takie, że $x\in B\subseteq U$.
\end{enumerate}
\end{stwierdzenie}
\begin{proof}
Implikacja (i)$\Rightarrow$(ii) jest definicją bazy.
W drugą stronę $U$ jest sumą wszystkich elementów
$\mathcal B$ zawartych w $U$, a więc zbiorem otwartym.
Gdy $U=\varnothing$, warunek (ii) jest spełniony,
a $U$ jest otwarty z aksjomatu (T1).
\end{proof}

\begin{przyklad}[Kule metryczne jako baza]
\label{ex:topog-kule-baza}
W dowolnej przestrzeni metrycznej $(X,d)$ rodzina
$\mathcal B_d=\{B_d(x,r):x\in X,\ r>0\}$
jest bazą topologii metrycznej. Istotnie, z
Definicji~\ref{def:topog-otwarty-metrycznie} wynika,
że dla każdego otwartego $U\subseteq X$ i $x\in U$
istnieje $B_d(x,r)$ zawarta w $U$. Każda kula należy
przy tym do topologii metrycznej na mocy
Stwierdzenia~\ref{prop:topog-kule-otwarte-sfery-domkniete}.
Są to dokładnie warunki z Definicji~\ref{def:topog-baza}.
Jeżeli $X=\varnothing$, baza jest pusta i warunek
jej definicji zachodzi bez konieczności wyboru punktu.
\end{przyklad}

\begin{przyklad}[Baza o końcach wymiernych]
\label{ex:topog-wymierna-baza}
W zwykłej topologii na $\R$ rodzina
$\mathcal B_{\mathbb Q}=\{(p,q):p,q\in\Q,\ p<q\}$
jest bazą. Niech $U$ będzie otwarty i $x\in U$.
Istnieje $\varepsilon>0$ takie, że
$(x-\varepsilon,x+\varepsilon)\subseteq U$.
Z gęstości $\Q$ w $\R$ wybieramy liczby wymierne
$p,q$ tak, aby $x-\varepsilon<p<x<q<x+\varepsilon$.
Wtedy $x\in(p,q)\subseteq U$. Rodzina jest przeliczalna: jest obrazem przeliczalnego
zbioru $\{(p,q)\in\Q\times\Q:p<q\}$ przy odwzorowaniu,
które parze liczb $(p,q)$ przyporządkowuje przedział $(p,q)$.
Własność tę wykorzystamy przy aksjomatach przeliczalności.
\end{przyklad}

\begin{przyklad}[Topologia dolnej granicy]
\label{ex:topog-sorgenfrey}
Niech $X=\R$ i
$\mathcal B_{\mathrm S}=\{[a,b):a,b\in\R,\ a<b\}$.
Rodzina pokrywa $\R$, gdyż $x\in[x,x+1)$.
Jeżeli $x\in[a,b)\cap[c,d)$, to
$x<\min\{b,d\}$ i
$[x,\min\{b,d\})\subseteq[a,b)\cap[c,d)$.
Kryterium bazy wyznacza topologię
$\tau_{\mathrm{So}}$, zwaną \emph{topologią dolnej granicy}
(topologią prostej Sorgenfreya).
Każdy zbiór otwarty w zwykłej topologii na $\R$
jest otwarty w $\tau_{\mathrm{So}}$: dla $x\in U$
można wybrać $\varepsilon>0$ z
$(x-\varepsilon,x+\varepsilon)\subseteq U$, skąd
$x\in[x,x+\varepsilon)\subseteq U$.
Natomiast $[0,1)$ jest otwarty w $\tau_{\mathrm{So}}$,
ale nie jest otwarty w topologii zwykłej,
bo żadna kula wokół zera nie mieści się w $[0,1)$.
Topologia dolnej granicy jest zatem ściśle silniejsza
od topologii zwykłej.
\end{przyklad}

\section{Podbazy i topologie generowane}
\label{sec:topog-podbazy}

\begin{definicja}[Topologia generowana przez rodzinę zbiorów]
\label{def:topog-topologia-generowana}
Niech $X$ będzie zbiorem, a $\mathcal S\subseteq\mathcal P(X)$.
\emph{Topologią generowaną przez $\mathcal S$} nazywamy
najmniejszą względem zawierania topologię na $X$,
która zawiera wszystkie elementy $\mathcal S$.
Rodzinę $\mathcal S$ nazywamy \emph{podbazą} tej topologii.
\end{definicja}

Istnienie najmniejszej topologii nie jest założeniem:
wynika z następnego twierdzenia. W definicji
\emph{nie} zakładamy, że $\mathcal S$ pokrywa $X$.

\begin{twierdzenie}[Konstrukcja topologii z podbazy]
\label{thm:topog-podbaza-konstrukcja}
Niech $X$ będzie dowolnym zbiorem i
$\mathcal S\subseteq\mathcal P(X)$. Niech $\mathcal B$
będzie rodziną wszystkich przecięć skończonych podrodzin
$\mathcal S$, przy czym przecięcie pustej podrodziny
jest równe $X$. Wówczas $\mathcal B$ jest bazą pewnej
topologii, a topologia ta jest dokładnie topologią
generowaną przez $\mathcal S$. Każdy jej zbiór otwarty
jest sumą skończonych przecięć elementów $\mathcal S$.
\end{twierdzenie}
\begin{proof}
Ponieważ $X\in\mathcal B$, każdy punkt jest zawarty
w pewnym elemencie $\mathcal B$. Przecięcie dwóch
skończonych przecięć elementów $\mathcal S$ jest
znów takim skończonym przecięciem. Zatem $\mathcal B$
spełnia (B1)--(B2) z
Twierdzenia~\ref{thm:topog-kryterium-bazy}.
Oznaczmy wyznaczoną topologię przez $\tau$.
Każdy $S\in\mathcal S$ należy do $\mathcal B$ jako
przecięcie jednoelementowej rodziny, zatem $\mathcal S
\subseteq\tau$.

Niech $\rho$ będzie dowolną topologią na $X$
zawierającą $\mathcal S$. Z (T1) i (T3) wynika,
że każde skończone przecięcie elementów $\mathcal S$
należy do $\rho$, a z (T2), że należą do niej także
wszystkie sumy takich przecięć. Zatem $\tau\subseteq\rho$.
Dowodzi to minimalności i jednoznaczności $\tau$.
\end{proof}

\begin{przyklad}[Podbaza nie musi być bazą]
\label{ex:topog-podbaza-niebaza}
Niech $X=\{a,b,c\}$ i
$\mathcal S=\{\{a,b\},\{b,c\}\}$.
Przecięcie obu elementów $\mathcal S$ jest równe
$\{b\}$, ale w $\mathcal S$ nie ma zbioru
zawierającego $b$ i zawartego w $\{b\}$.
Zatem $\mathcal S$ nie spełnia (B2), choć pokrywa $X$.
Topologia przez nią generowana to
\[
\tau=\{\varnothing,\{b\},\{a,b\},\{b,c\},X\}.
\]
Istotnie, przecięcia skończone dają tylko $X$,
$\{a,b\}$, $\{b,c\}$ i $\{b\}$,
a sumy tych zbiorów, wraz z sumą pustą,
dają dokładnie wypisaną rodzinę.
\end{przyklad}

\begin{przyklad}[Pusta podbaza]
\label{ex:topog-pusta-podbaza}
Gdy $\mathcal S=\varnothing$, jedynym skończonym
przecięciem zbiorów z $\mathcal S$ jest $X$
(przecięcie puste). Sumy podrodzin $\{X\}$
dają $\varnothing$ oraz $X$, a więc topologię
antydyskretną. Jest to zgodne z tym, że żadnych
zbiorów nie nakazano uczynić otwartymi.
\end{przyklad}

\section{Porównywanie topologii na ustalonym zbiorze}
\label{sec:topog-porownywanie}

\begin{definicja}[Topologia słabsza i silniejsza]
\label{def:topog-silniejsza}
Niech $\tau_1$ i $\tau_2$ będą topologiami na tym
samym zbiorze $X$. Mówimy, że $\tau_1$ jest
\emph{słabsza} (albo \emph{grubsza}) od $\tau_2$,
a $\tau_2$ jest \emph{silniejsza} (albo \emph{drobniejsza})
od $\tau_1$, jeżeli $\tau_1\subseteq\tau_2$.
Jeśli zawieranie jest właściwe, mówimy odpowiednio
o topologii \emph{ściśle słabszej} i \emph{ściśle silniejszej}.
\end{definicja}

To porównanie dotyczy liczby i układu zbiorów otwartych,
a nie wartości odległości. Topologia dyskretna jest
najsilniejsza, a antydyskretna najsłabsza ze wszystkich
topologii na $X$.

\begin{stwierdzenie}[Kryterium porównania przez bazy]
\label{prop:topog-porownanie-bazy}
Niech $\tau_1,\tau_2$ będą topologiami na $X$,
a $\mathcal B_1$ bazą $\tau_1$.
Wówczas $\tau_1\subseteq\tau_2$ wtedy i tylko wtedy,
gdy każdy zbiór $B\in\mathcal B_1$ należy do $\tau_2$.
Jeżeli $\mathcal B_2$ jest bazą $\tau_2$, jest to
równoważne warunkowi: dla wszystkich $B\in\mathcal B_1$
i $x\in B$ istnieje $C\in\mathcal B_2$ takie,
że $x\in C\subseteq B$.
\end{stwierdzenie}
\begin{proof}
Jeżeli $\tau_1\subseteq\tau_2$, każdy element
$\mathcal B_1\subseteq\tau_1$ należy do $\tau_2$.
Odwrotnie, jeśli wszystkie elementy $\mathcal B_1$
są otwarte w $\tau_2$, to dowolny $U\in\tau_1$
jako ich suma również należy do $\tau_2$.
Ostatnia równoważność wynika z kryterium otwartości
przez bazę $\mathcal B_2$.
\end{proof}

\begin{przyklad}[Topologie nieporównywalne]
\label{ex:topog-nieporownywalne}
Na $X=\{a,b\}$ określmy
$\tau_a=\{\varnothing,\{a\},X\}$ i
$\tau_b=\{\varnothing,\{b\},X\}$.
Obie rodziny są topologiami, ale
$\{a\}\in\tau_a\setminus\tau_b$ i
$\{b\}\in\tau_b\setminus\tau_a$.
Żadna z tych topologii nie jest słabsza od drugiej.
Na zbiorze dwuelementowym istnieją dokładnie cztery
topologie: antydyskretna, dyskretna oraz $\tau_a,\tau_b$.
Wynika to z tego, że $\varnothing$ i $X$ są obowiązkowe,
a każdy z dwóch pozostałych podzbiorów można dołączyć
lub pominąć, nie naruszając aksjomatów.
\end{przyklad}

\begin{figure}[ht]
\centering
\begin{tikzpicture}[>=Latex, every node/.style={font=\small},
 box/.style={draw,rounded corners,align=center,inner sep=4pt}]
\node[box] (top) at (0,2.6) {$\mathcal P(X)$};
\node[box] (left) at (-2.5,1.3) {$\tau_a$};
\node[box] (right) at (2.5,1.3) {$\tau_b$};
\node[box] (bottom) at (0,0) {$\{\varnothing,X\}$};
\draw (bottom)--(left)--(top)--(right)--(bottom);
\end{tikzpicture}
\caption{Cztery topologie na $X=\{a,b\}$ uporządkowane
przez zawieranie; kierunek ku górze oznacza topologię silniejszą.}
\label{fig:topog-cztery-topologie}
\end{figure}

\begin{twierdzenie}[Przecięcia topologii i topologia generowana]
\label{thm:topog-przeciecia-topologii}
Niech $X$ będzie zbiorem i $(\tau_i)_{i\in I}$
dowolną rodziną topologii na $X$. Wówczas
$\bigcap_{i\in I}\tau_i$ jest topologią na $X$,
przy czym dla $I=\varnothing$ przyjmujemy
$\bigcap_{i\in\varnothing}\tau_i=\mathcal P(X)$.
W szczególności dla każdego
$\mathcal S\subseteq\mathcal P(X)$ topologia
generowana przez $\mathcal S$ jest przecięciem
wszystkich topologii na $X$ zawierających $\mathcal S$.
\end{twierdzenie}
\begin{proof}
Dla $I\neq\varnothing$ zbiory $\varnothing,X$
należą do każdej $\tau_i$, więc także do ich przecięcia.
Jeżeli $(U_j)_{j\in J}$ jest rodziną zbiorów należących
do wszystkich $\tau_i$, to dla każdego $i\in I$
zbiór $\bigcup_{j\in J}U_j$ należy do $\tau_i$;
zatem należy do ich przecięcia. Analogicznie,
jeżeli $U,V$ należą do wszystkich $\tau_i$, to
$U\cap V$ należy do wszystkich $\tau_i$.
Dla $I=\varnothing$ otrzymujemy topologię dyskretną.

Rodzina wszystkich topologii zawierających $\mathcal S$
nie jest pusta, ponieważ zawiera $\mathcal P(X)$.
Jej przecięcie jest topologią zawierającą $\mathcal S$
i zawartą w każdej innej topologii o tej własności.
Zatem jest to najmniejsza topologia zawierająca
$\mathcal S$, zgodnie z Definicją~\ref{def:topog-topologia-generowana}.
\end{proof}

\begin{przyklad}[Porównanie trzech topologii na prostej]
\label{ex:topog-porownanie-r}
Na $\R$ topologia koskończona jest ściśle słabsza od
zwykłej topologii euklidesowej. Istotnie, dla każdego
$a\in\R$ zbiór $\R\setminus\{a\}$ jest otwarty:
jeśli $x\ne a$, to
$B_d(x,|x-a|/2)\subseteq\R\setminus\{a\}$.
Dopełnienie dowolnego zbioru skończonego jest zatem
przecięciem skończenie wielu takich zbiorów otwartych,
a więc jest otwarte. Jednak $(-1,1)$ nie jest otwarty
w topologii koskończonej, ponieważ jego dopełnienie
jest nieskończone.

Topologia euklidesowa i topologia koprzeliczalna
na $\R$ są nieporównywalne. Zbiór
$\R\setminus\Q$ jest koprzeliczalnie otwarty,
ale nie jest euklidesowo otwarty: każdy przedział
wokół liczby niewymiernej zawiera liczbę wymierną.
Z kolei $(-1,1)$ jest euklidesowo otwarty,
ale jego dopełnienie nie jest przeliczalne,
więc nie jest koprzeliczalnie otwarty.
\end{przyklad}

\section{Podsumowanie i dalsze powiązania}
\label{sec:topog-podsumowanie}

W przestrzeni metrycznej odległość umożliwia definiowanie
kul i sfer; z kul wyprowadzamy pojęcie zbioru otwartego.
Przestrzeń topologiczna to zbiór wyposażony w rodzinę zbiorów
otwartych, zamkniętą na dowolne sumy i skończone przecięcia.
Te same informacje można opisywać za pomocą zbiorów
domkniętych, otoczeń, baz albo podbaz. Znajomość bazy
pozwala sprawdzać otwartość lokalnie, w otoczeniu każdego
punktu. Podbaza daje z kolei sposób wygenerowania
najmniejszej topologii zawierającej zadaną rodzinę zbiorów.

W następnym rozdziale przejdziemy do operacji na podzbiorach
przestrzeni topologicznej: wnętrza, domknięcia, brzegu i zbioru
pochodnego. Następnie zbadamy ciągłość odwzorowań, dla której
istotna będzie właśnie rodzina zbiorów otwartych,
a nie wybrana metryka.
